%
% 32-laurent.tex -- Die Laurent-Reihe
%
% (c) 2026 Prof Dr Andreas Müller
%
\subsection{Die Laurent-Reihe}
Auf einem Kreisgebiet lässt sich eine holomorphe Funktion immer in eine
Potenzreihe entwickeln.
Auf einem Ringgebiet ist dies nicht mehr möglich.
An die Stelle der Potenzreihe tritt ein sogenannte Laurent-Reihe.

\begin{definition}[Laurent-Reihe]
\label{buch:analytisch:laurent:def:laurentreihe}
Eine \emph{Laurent-Reihe} ist eine Reihe der Form
\index{Laurent-Reihe}%
\begin{equation}
f(a)
=
\sum_{k=0}^\infty a_k(a-z_0)^k
+
\sum_{k=1}^\infty b_k(a-z_0)^{-k}
\label{buch:analytisch:laurent:def:laurent}
\end{equation}
mit $a_k\in\mathbb{C}$, $k\in\mathbb{N}$ und $b_k\in\mathbb{C}$,
$k\in\mathbb{N}\setminus\{0\}$.
\end{definition}

Die erste Reihe in
\eqref{buch:analytisch:laurent:def:laurent}
ist eine gewöhnliche Potenzreihen, deren Konvergenzradius wir mit $r_2$
bezeichnen.
Die zweite Reihe können wir auch als
\begin{equation}
\sum_{k=1}^\infty b_k(a-z_0)^{-k}
=
\sum_{k=1}^\infty b_k\biggl(\frac{1}{a-z_0}\biggr)^{k}
\label{buch:analytisch:laurent:def:laurent2}
\end{equation}
schreiben, also als Potenzreihe in $1/(a-z_0)$.
Der Konvergenzradius $\varrho$ dieser Potenzreihe kann mit der
Cauchy-Hadamard-Formel berechnet werden.
Die Reihe~\eqref{buch:analytisch:laurent:def:laurent2}
konvergiert dann für $|a-z_0| > 1/\varrho = r_1$.
Nach dem Satz von Cauchy-Hadamard ist die Laurent-Reihe ist also 
auf dem Ringgebiet $R(z_0;r_1,r_2)$ mit den Radien
\begin{align*}
\frac{1}{r_1}
&=
\limsup_{n\to\infty}\!\sqrt[n]{|b_n|}
&&\text{und}&
r_2
&=
\limsup_{n\to\infty}\!\sqrt[n]{|a_n|}.
\end{align*}
und $r_1 < r_2$ konvergent.

\begin{satz}[Laurent-Reihe]
\label{buch:analytisch:laurent:satz:laurentreihe}
Ist $f$ eine auf einer Umgebung des Ringgebietes $R(z_0;r_1,r_2)$ definierte
holomorphe Funktion, dann kann sie in eine Laurent-Reihe entwickelt werden.
Ist $\gamma$ ein beliebiger Weg im Ringgebiet, der $z_0$ genau einmal
im Gegenuhrzeigersinn umläuft, dann sind die 
Koeffizienten durch
\begin{align}
a_k
&=
\frac{1}{2\pi i}
\oint_{\gamma} 
\frac{f(z)}{(z-z_0)^{k+1}}
\,dz
&&\text{und}&
b_k
&=
\frac{1}{2\pi i}
\oint_{\gamma} 
f(z)(z-z_0)^{k-1}
\,dz
\label{buch:analytisch:laurent:eqn:laurentkoeffizienten}
\end{align}
gegeben.
\end{satz}

\begin{proof}
Wir berechnen jetzt die beiden Integrale in
\eqref{buch:analytisch:laurent:eqn:zweiintegrale} separat mit
verschiedenen Entwicklungen des Integranden.
Das erste Integral über den äusseren Randkreis $\gamma_2$ des Ringgebietes
lässt sich wie früher in eine geometrische Reihe
\begin{equation}
\frac{1}{2\pi i}
\oint_{\gamma_2}
\frac{f(z)}{z-a}
\,dz
=
\sum_{k=0}^\infty
\biggl(
\underbrace{
\frac{1}{2\pi i}
\oint_{\gamma_2}
\frac{f(z)}{(z-z_0)^{k+1}}
\,dz
}_{\displaystyle =a_k}
\biggr)
(a-z_0)^k
=
\sum_{k=0}^\infty a_k(a-z_0)^k
\label{buch:analytisch:laurent:eqn:integral1}
\end{equation}
entwickeln.

Für das Integral über den inneren Randkreis  $\gamma_1$
verwenden wir die geometrische Reihe
\begin{align*}
\frac{1}{z-a\mathstrut}
&=
\frac{1}{z-z_0-(a-z_0)\mathstrut}
\\
&=
-
\frac{1}{a-z_0}\cdot \frac{1}{1-\displaystyle\frac{z-z_0\mathstrut}{a-z_0}}
=
-
\frac{1}{a-z_0}
\sum_{k=0}^\infty \biggl(\frac{z-z_0}{a-z_0}\biggr)^k
\\
&=
-
\sum_{k=0}^\infty \frac{1}{(a-z_0)^{k+1}}(z-z_0)^k.
\end{align*}
Damit kann das Integral über $\gamma_1$ als
\begin{align}
-
\frac{1}{2\pi i}
\oint_{\gamma_1}
\frac{f(z)}{z-a}
\,dz
&=
\frac{1}{2\pi i}
\oint_{\gamma_1}
\sum_{k=0}^\infty
f(z)(z-z_0)^k
\frac{1}{(a-z_0)^{k+1}}
\,dz
\notag
\\
&=
\sum_{k=0}^\infty
\biggl(
\frac{1}{2\pi i}
\oint_{\gamma_1}
f(z) (z-z_0)^k
\,dz
\biggr)
\cdot
\frac{1}{(a-z_0)^{k+1}}.
\notag
\intertext{geschrieben werden.
Durch Verschiebung des Summationsindex wird daraus}
-
\frac{1}{2\pi i}
\oint_{\gamma_1}
\frac{f(z)}{z-a}
\,dz
&=
\sum_{k=1}^\infty
\biggl(
\underbrace{
\frac{1}{2\pi i}
\oint_{\gamma_1}
f(z) (z-z_0)^{k-1}
\,dz
}_{\displaystyle =b_k}
\biggr)
\cdot
\frac{1}{(a-z_0)^k}
\notag
\\
&=
\sum_{k=1}^\infty
b_k
\cdot
\frac{1}{(a-z_0)^k}.
\label{buch:analytisch:laurent:eqn:integral2}
\end{align}

Die beiden Reihenentwicklungen
\eqref{buch:analytisch:laurent:eqn:integral1}
und
\eqref{buch:analytisch:laurent:eqn:integral2}
können wir jetzt zu
\begin{align*}
f(a)
&=
\underbrace{
\sum_{k=0}^\infty
a_k
\cdot 
(a-z_0)^k
}_{\displaystyle\text{\eqref{buch:analytisch:laurent:eqn:integral1}}}
+
\underbrace{
\sum_{k=1}^\infty
b_k
\cdot 
(a-z_0)^{-k}
}_{\displaystyle\text{\eqref{buch:analytisch:laurent:eqn:integral1}}}
\end{align*}
zusammensetzen, dies ist die versprochene Laurent-Reihe.

Jetzt muss noch gezeigt werden, dass die Integrale über $\gamma_1$
bzw.~$\gamma_2$ durch Integrale über den Weg $\gamma$ irgendwo im
Ringgebiet ersetzt werden können.
Die Integranden $f(z)(z-z_0)^k$, $k\in\mathbb{Z}$, der Integrale
sind im ganzen Ringgebiet holomorph, da $f(z)$ dort holomorph ist
und die einzige mögliche Polstelle von $(z-z_0)^k$ im Auge des
Ringgebietes liegt.
Die Deformation des Weges $\gamma_1$ in den Weg $\gamma$ oder des
Weges $\gamma_2$ in den Weg $\gamma$ ändert daher das Integral
nicht.
\end{proof}
